\begin{thebibliography}{99}
\small
\bibitem{kim2024} J.~Kim. Recording and describing poker hands. In \emph{IEEE Conference on Games (CoG)}, 2024. arXiv:2312.11753.
\bibitem{kim2025} J.~Kim. A dataset of poker hand histories (v3). Zenodo, 2025. doi:10.5281/zenodo.10796885. Mirrored as \texttt{takara-ai/poker\_hands} on Hugging Face (CC BY 4.0).
\bibitem{brown2019} N.~Brown and T.~Sandholm. Superhuman AI for multiplayer poker. \emph{Science}, 365(6456):885--890, 2019.
\bibitem{rabiner1989} L.~R.~Rabiner. A tutorial on hidden Markov models and selected applications in speech recognition. \emph{Proceedings of the IEEE}, 77(2):257--286, 1989.
\bibitem{dempster1977} A.~P.~Dempster, N.~M.~Laird and D.~B.~Rubin. Maximum likelihood from incomplete data via the EM algorithm. \emph{Journal of the Royal Statistical Society B}, 39(1):1--38, 1977.
\bibitem{bengio1995} Y.~Bengio and P.~Frasconi. An input output HMM architecture. In \emph{Advances in Neural Information Processing Systems 7}, 1995.
\bibitem{smith2009} G.~Smith, M.~Levere and R.~Kurtzman. Poker player behavior after big wins and big losses. \emph{Management Science}, 55(9):1547--1555, 2009.
\bibitem{eil2014} D.~Eil and J.~W.~Lien. Staying ahead and getting even: Risk attitudes of experienced poker players. \emph{Games and Economic Behavior}, 87:50--69, 2014.
\bibitem{thaler1990} R.~H.~Thaler and E.~J.~Johnson. Gambling with the house money and trying to break even: The effects of prior outcomes on risky choice. \emph{Management Science}, 36(6):643--660, 1990.
\bibitem{hmmlearn} hmmlearn developers. hmmlearn: hidden Markov models in Python (version 0.3.3). \url{https://github.com/hmmlearn/hmmlearn}.
\bibitem{numpy} C.~R.~Harris et al. Array programming with NumPy. \emph{Nature}, 585:357--362, 2020.
\end{thebibliography}
